\documentclass[aspectratio=169,10pt,spanish]{beamer}

\input{../../../_preambulo_beamer.tex}
\input{../../../_comandos_beamer.tex}

\selectlanguage{spanish}

\title{MA1001B - Modelación Estadística para la Toma de Decisiones}
\subtitle{Lección 14: Ensayos de Bernoulli y el modelo binomial}
\author{}
\institute{Tecnol\'ogico de Monterrey}
\date{}

\begin{document}

\begin{frame}[plain]
  \vspace{1.2cm}
  {\Large\bfseries MA1001B - Modelación Estadística para la Toma de Decisiones\par}
  \vspace{0.7cm}
  {\large Lección 14: Ensayos de Bernoulli y el modelo binomial\par}
  \vspace{0.7cm}
  {\color{TecRojo}\rule{\textwidth}{1.2pt}}
  \vspace{0.8cm}

  Facultad MA1001B\\[0.3cm]
  Periodo\\[0.3cm]
  Tecnol\'ogico de Monterrey
\end{frame}

\begin{frame}{De un ensayo a un conteo}
  \begin{alertblock}{Pregunta guía}
    Si cada una de $n=20$ inspecciones es defectuosa con chance $p=0.08$,
    ¿cuál es $P(X=0)$, y qué se rompe si $p$ no es el mismo en cada lote?
  \end{alertblock}

  \vspace{0.6em}
  Al final de esta lección, usted puede:
  \begin{itemize}
    \item enunciar el ensayo de Bernoulli y los supuestos binomiales;
    \item calcular $P(X=0)$ y $P(X\le 2)$;
    \item reportar $np$ y $np(1-p)$;
    \item explicar por qué tipos de lote mezclados no son binomiales.
  \end{itemize}

  \vfill
  \scriptsize Todas las inspecciones usadas hoy son sintéticas. No se usan datos reales de empresa.
\end{frame}

\begin{frame}{Ensayo de Bernoulli, luego $n$ copias independientes}
  \small
  Una inspección: $P(\text{defectuoso})=p=0.08$, $P(\text{limpio})=q=0.92$.
  Repita $n=20$ veces, de forma independiente, con el mismo $p$. Entonces
  $X\sim\mathrm{Binomial}(n=20,p=0.08)$ y
  \[
    E[X]=np=1.6,\qquad \Var(X)=np(1-p)=1.472.
  \]
\end{frame}

\begin{frame}[fragile]{Paso 1: Ensayos de Bernoulli}
  \begin{columns}[T]
    \begin{column}{0.42\textwidth}
      \small
      \begin{lstlisting}
p = 0.08
q = 1 - p
n = 20
mean = n * p
      \end{lstlisting}

      \begin{block}{Salida verificada}
        $p=0.080000$\\
        $q=0.920000$\\
        $n=20$\\
        $np=1.600000$
      \end{block}
    \end{column}
    \begin{column}{0.58\textwidth}
      \centering
      \includegraphics[width=\textwidth]{../figuras/binomial_01_ensayos_bernoulli.png}
      \vspace{0.2em}
      \scriptsize Una inspección de Bernoulli del Paso 1
    \end{column}
  \end{columns}
\end{frame}

\begin{frame}{Cola izquierda binomial}
  \small
  \[
    P(X=0)=(0.92)^{20}=0.188693,
  \]
  \[
    P(X\le 2)=0.787946.
  \]
  Estos valores usan \texttt{scipy.stats.binom} con $n=20$ y $p=0.08$.
\end{frame}

\begin{frame}[fragile]{Paso 2: FMP, FDA, momentos}
  \begin{columns}[T]
    \begin{column}{0.40\textwidth}
      \small
      \begin{lstlisting}
model = stats.binom(n=20, p=0.08)
p0 = model.pmf(0)
p2 = model.cdf(2)
      \end{lstlisting}

      \begin{block}{Salida verificada}
        $P(X=0)=0.188693$\\
        $P(X\le 2)=0.787946$\\
        $\Var(X)=1.472000$
      \end{block}
    \end{column}
    \begin{column}{0.60\textwidth}
      \centering
      \includegraphics[width=\textwidth]{../figuras/binomial_02_fmp_fda.png}
      \vspace{0.2em}
      \scriptsize FMP binomial del Paso 2
    \end{column}
  \end{columns}
\end{frame}

\begin{frame}{Paso 3: $p$ no constante entre lotes}
  \begin{columns}[T]
    \begin{column}{0.42\textwidth}
      \small
      La mitad de los lotes tiene $p=0.02$ y la mitad tiene $p=0.14$. El $p$
      promedio sigue siendo $0.08$, y $E[X]$ sigue siendo $1.600000$.

      \vspace{0.4em}
      $P(X=0)$ de la mezcla $=0.358291$ versus binomial $0.188693$. Varianza
      de la mezcla $2.840000$ versus $1.472000$.

      \vspace{0.3em}
      \textbf{Límite:} $p$ no constante entre lotes rompe el binomial. La
      Lección 15 compara dos políticas binomiales con la misma media.
    \end{column}
    \begin{column}{0.58\textwidth}
      \centering
      \includegraphics[width=\textwidth]{../figuras/binomial_03_p_no_constante.png}
      \vspace{0.2em}
      \scriptsize Mezcla versus binomial del Paso 3
    \end{column}
  \end{columns}
\end{frame}

\begin{frame}{Verificación de conceptos}
  \begin{enumerate}
    \item ¿Cuáles son los dos resultados de una inspección de Bernoulli aquí?
    \item Calcule $np$ para $n=20$ y $p=0.08$.
    \item ¿Por qué $P(X=0)=(0.92)^{20}$?
    \item ¿Por qué promediar $p=0.02$ y $p=0.14$ no restaura el binomial?
  \end{enumerate}

  \vspace{0.6em}
  \begin{block}{Conexión con la Lección 15}
    La sesión puede continuar directamente con dos políticas binomiales que
    comparten $np$ pero no el mismo riesgo de cola.
  \end{block}

  \vfill
  \scriptsize Fuente: OpenStax, \textit{Introductory Business Statistics 2e},
  Sección 4.2. CC BY-NC-SA.
\end{frame}

\appendix



\end{document}
